\documentclass[addpoints]{exam}
% \documentclass[addpoints,12pt]{exam}

\usepackage[slovene]{babel}
\usepackage{amsfonts}
\usepackage{amssymb}
\usepackage{amsmath}
\usepackage{float}
\usepackage{graphicx}
\usepackage{enumitem}
\usepackage{color}
\usepackage[gen]{eurosym}
\usepackage{newunicodechar}
\newunicodechar{€}{\euro}

%Komentar naslednje vrstice glede na verzijo testa -- rešitve ali prostor za odgovore:
% \printanswers

% \linespread{2}


\pointsinrightmargin
\marginbonuspointname{}


\renewcommand{\solutiontitle}{\noindent\textbf{Rešitve in točkovanje:}\par\noindent}

\colorgrids
\definecolor{GridColor}{gray}{0.7}
\setlength{\gridsize}{7mm}

\extrawidth{5mm}
\setlength{\rightpointsmargin}{1.5cm}

\begin{document}

\runningfooter{}{}{}

\begin{center}
    \fbox{\fbox{\parbox{15cm}{\centering
    Popravljanje in pridobivanje ocen -- konferenca \\ 1. letnik -- test D \\ 15. junij 2026 \\ (Racionalna števila; Realna števila; Koordinatni sistem v ravnini; Funkcija; Premica)}}}
\end{center}

\vspace{0.1in}
\makebox[\textwidth]{Ime in priimek, oddelek:\enspace\hrulefill}


\begin{center}
    \hqword{Naloga}
    \hpword{Možne točke}
    \chbpword{Dodatne točke}
    \hsword{Dosežene točke}
    \htword{Skupaj}  
    % \combinedpointtable[h][questions]
    \gradetable[h][questions]

\end{center}

\vspace{0.02in}


\begin{questions}

    %%%%% racionalna števila
    
    % 1. vprašanje

    \question[4]
    Izračunajte natančno vrednost izraza.
    % \begin{parts}
    %     \part[4]
        $$\dfrac{35\cdot \frac{1}{7}+\frac{7}{8}}{2\frac{1}{3}+\frac{7}{4}}-4\frac{8}{9}:\frac{11}{3}$$

        \begin{solution}[9cm]
            \begin{align*} 
                \dfrac{\frac{2}{3}+1\frac{5}{6}\cdot 3}{\frac{2}{7}:\frac{4}{3}-1\frac{1}{14}+2}&=\dfrac{\frac{2}{3}+\frac{11}{6}\cdot 3}{\frac{2}{7}\cdot\frac{3}{4}-\frac{15}{14}+\frac{28}{14}}= \\
                                                                                                &=\dfrac{\frac{2}{3}+\frac{11}{2}}{\frac{1}{7}\cdot\frac{3}{2}+\frac{13}{14}}= \\
                                                                                                &=\dfrac{\frac{4}{6}+\frac{33}{6}}{\frac{3}{14}+\frac{13}{14}}= \\
                                                                                                &=\dfrac{\frac{37}{6}}{\frac{16}{14}}= \\
                                                                                                &=\dfrac{37\cdot 14}{6\cdot 16}= \\
                                                                                                &=\dfrac{259}{48} 
            \end{align*}
            Za pravilno seštevanje/odštevanje ulomkov $1$ točka; pravilno množenje ulomkov $1$ točka; pravilno razrešen dvojni ulomek $1$ točka, okrajšan rezultat $1$ točka.
        \end{solution}




          \question[4]
        Ceno puloverja so zvišali za $20~\%$, a ker ni šel v prodajo, so ga nato pocenili za $30~\%$.
        Po drugi pocenitvi ga je kupil Matej in zanj plačal $30.24~\euro$. Določite prvotno ceno puloverja.


        \begin{solution}[4cm]
            \begin{align*}
                0.70\cdot \left(1.10\cdot x\right) &=115.5~\euro \\
                                     0.77 \cdot x &= 115.5~\euro \\
                                                x &= \dfrac{115.5}{0.77}~\euro \\
                                                x &= 150~\euro
            \end{align*}
            Pravilen zapis in reševanje enačbe $1$ točka, pravilen rezultat $1$ točka, zapis pravilnega odgovora $1$ točka.
        \end{solution}



        \newpage
      % 3. vprašanje

        % \question[5]
        % Poenostavite izraz.
    
        %     $$\dfrac{1}{x+2}-\dfrac{3}{x-2}-4\left(4-x^2\right)^{-1}$$
    
        %     \begin{solution}[8cm]
        %         \begin{align*}
        %             \left(\dfrac{6-x}{x^2+6x}-\dfrac{x}{36-x^2}\right)\left(\dfrac{2x-6}{x^2+6x}\right)^{-1}+\dfrac{x}{6-x} &= \left(\dfrac{6-x}{x(x+6)}-\dfrac{x}{(6-x)(6+x)}\right)\dfrac{x(x+6)}{2(x-3)}+\dfrac{x}{6-x}= \\
        %                                                                                                                     &= \left(\dfrac{(6-x)^2}{x(x+6)(6-x)}-\dfrac{x^2}{x(6-x)(6+x)}\right)\dfrac{x(x+6)}{2(x-3)}+\dfrac{x}{6-x}= \\
        %                                                                                                                     &= \dfrac{36-12x+x^2-x^2}{x(x+6)(6-x)}\dfrac{x(x+6)}{2(x-3)}+\dfrac{x}{6-x}= \\
        %                                                                                                                     &= \dfrac{12(3-x)}{x(x+6)(6-x)}\dfrac{x(x+6)}{2(x-3)}+\dfrac{x}{6-x}= \\
        %                                                                                                                     &= \dfrac{-6}{6-x}+\dfrac{x}{6-x}= \\
        %                                                                                                                     &= \dfrac{-6+x}{6-x}= \\
        %                                                                                                                     &= -1
        %         \end{align*}
        %         Za pravilno faktorizacijo izrazov $1$ točka, pravilno razčlenjevanje izrazov $1$ točka, pravilno seštevanje ulomkov $1$ točka, pravilna obratna vrednost ulomka $1$ točka, pravilno množenje ulomkov $1$ točka, pravilno krajšanje ulomkov $1$ točka, okrajšan rezultat $1$ točka.
        %     \end{solution}
    
    
   
   
   
   


%%%% realna števila


    \question[4]
    Izračunajte natančno vrednost izraza.
    $$ \sqrt{6\cdot 8\cdot 12}-\left(\sqrt{3}-1\right)^2+\sqrt{3}\left(2+\sqrt{3}\right)^{-1} $$

    \begin{solution}[9cm]
        \begin{align*} 
            \dfrac{\sqrt{5}}{\sqrt{5}+1}-\dfrac{5-3\sqrt{5}}{4\sqrt{5}}+\left(3+\sqrt{5}\right)^2-\sqrt{125}&=\dfrac{\sqrt{5}(\sqrt{5}-1)}{(\sqrt{5}+1)(\sqrt{5}-1)}-\dfrac{(5-3\sqrt{5})\sqrt{5}}{4\sqrt{5}\sqrt{5}}+9+6\sqrt{5}+5-5\sqrt{5}= \\
                                                                                                            &=\dfrac{5-\sqrt{5}}{5-1}-\dfrac{5\sqrt{5}-3\cdot 5}{4\cdot 5}+14+\sqrt{5}= \\
                                                                                                            &=\frac{5-\sqrt{5}}{4}-\dfrac{\sqrt{5}-3}{4}+14+\sqrt{5}= \\
                                                                                                            &=\frac{8-2\sqrt{5}}{4}+14+\sqrt{5}= \\
                                                                                                            &=16+\frac{1}{2}\sqrt{5}
        \end{align*}
        Za racionalizacijo $2$ točki, za izračun kvadrata dvočlenika $1$ točka, delno korenjenje $1$ točka, izračun $1$ točka, končne rezultat $1$ točka.
    \end{solution}
    
    

    
        

    
    % 3. vprašanje
    \question[5]
    Rešite sistem neenačb, rešitev grafično upodobite in zapišite z intervalom.
    $$ \left(3x-5>\frac{2}{3}x+2\right)\land\left(\frac{2}{5}x\geq x-3\right) $$
    
    
    \begin{solution}[8cm]
        Rešitev prve enačbe $x>2$ ali $x\in(2,\infty)$ $1$ točka, rešitev druge enačbe $x\geq 3$ ali $x\in[3,\infty)$ $1$ točka, rešitev sistema $x\geq 3$ ali $x\in[3,\infty)$ $1$ točka, zapis končne rešitve z intervalom $1$ točka, grafična upodobitev $1$ točka.
    \end{solution}


    



    \newpage

    % 7. vprašanje
    \question[5]
    Rešite enačbo in opišite njen geometrijski pomen ter rešitev grafično prikažite.
    $$  \left\lvert x+2\right\rvert=8 $$
    
    
    \begin{solution}[9cm]
        Opis geometrijskega pomena (oddaljenost števila $x$ od števila $-3$ je $2$) $1$ točka, grafična predstavitev $1$ točka, reševanje enačbe z upoštevanjem pogojev do $2$ točke, zapis rešitve $x\in\{-5,-1\}$ $1$ točka..
    \end{solution}









%%%% PKS, funkcija, premica

    % 1. vprašanje

    \question[4]
    Narišite dano množico točk v ravnini.
        $$[-1,2)\times \{-1,1,2\}$$

            \begin{solution}[7cm]
                ????
            \end{solution}

    
    
    \newpage

    % 2. vprašanje
    \question[4]
    Izračunajte ploščino in orientacijo trikotnika $\triangle ABC$ z oglišči v točkah $A(4,-3)$, $B(5,4)$ in $C(-4,1)$.
    
    
    \begin{solution}[5cm]
        ???
    \end{solution}

        
    

    % 5. vprašanje
    \question
    Dana je funkcija $$f(x)=\begin{cases}
    -2x-2, & x\leq 0 \\ -1, & 1<x<5
    \end{cases}.$$ 
    
    \begin{parts}
        \part[4]
        Narišite dano funkcijo.

            \begin{solution}[9cm]
                ???
            \end{solution}

        \part[1]
        Določite definicijsko območje funkcije $f$.
        
            \begin{solution}[1cm]
                ???
            \end{solution}

        \part[3]
        Izračunajte ničle in začetno vrednost funkcije $f$.

            \begin{solution}[2cm]
                ???
            \end{solution}


    \end{parts}
    
    





\end{questions}



\end{document}